% Options for packages loaded elsewhere
\PassOptionsToPackage{unicode}{hyperref}
\PassOptionsToPackage{hyphens}{url}
\PassOptionsToPackage{dvipsnames,svgnames,x11names}{xcolor}
%
\documentclass[
  13pt,
  a4paperpaper,
  12pt]{article}
\usepackage{amsmath,amssymb}
\usepackage{setspace}
\usepackage{iftex}
\ifPDFTeX
  \usepackage[T1]{fontenc}
  \usepackage[utf8]{inputenc}
  \usepackage{textcomp} % provide euro and other symbols
\else % if luatex or xetex
  \usepackage{unicode-math} % this also loads fontspec
  \defaultfontfeatures{Scale=MatchLowercase}
  \defaultfontfeatures[\rmfamily]{Ligatures=TeX,Scale=1}
\fi
\usepackage{lmodern}
\ifPDFTeX\else
  % xetex/luatex font selection
  \setmainfont[]{Times New Roman}
  \setmonofont[]{Courier New}
\fi
% Use upquote if available, for straight quotes in verbatim environments
\IfFileExists{upquote.sty}{\usepackage{upquote}}{}
\IfFileExists{microtype.sty}{% use microtype if available
  \usepackage[]{microtype}
  \UseMicrotypeSet[protrusion]{basicmath} % disable protrusion for tt fonts
}{}
\makeatletter
\@ifundefined{KOMAClassName}{% if non-KOMA class
  \IfFileExists{parskip.sty}{%
    \usepackage{parskip}
  }{% else
    \setlength{\parindent}{0pt}
    \setlength{\parskip}{6pt plus 2pt minus 1pt}}
}{% if KOMA class
  \KOMAoptions{parskip=half}}
\makeatother
\usepackage{xcolor}
\usepackage[margin=1.5cm,top=1.5cm,bottom=1.5cm,left=1.5cm,right=1.5cm,includeheadfoot]{geometry}
\usepackage{listings}
\newcommand{\passthrough}[1]{#1}
\lstset{defaultdialect=[5.3]Lua}
\lstset{defaultdialect=[x86masm]Assembler}
\setlength{\emergencystretch}{3em} % prevent overfull lines
\providecommand{\tightlist}{%
  \setlength{\itemsep}{0pt}\setlength{\parskip}{0pt}}
\setcounter{secnumdepth}{3}
% Minimal LaTeX preamble fragment for Pandoc-built documents
% Avoid \documentclass or loading hyperref/cleveref here (they are managed by the build pipeline)
% Provide safe unicode declarations and small custom macros only
\DeclareUnicodeCharacter{03BC}{\(\mu\)}
\DeclareUnicodeCharacter{03BB}{\(\lambda\)}
\DeclareUnicodeCharacter{03BD}{\(\nu\)}
\DeclareUnicodeCharacter{03C0}{\(\pi\)}
\DeclareUnicodeCharacter{03B5}{\(\epsilon\)}
\DeclareUnicodeCharacter{03B4}{\(\delta\)}

% Map common Unicode subscripts and degree symbol to LaTeX-friendly macros
\DeclareUnicodeCharacter{2080}{\textsubscript{0}}
\DeclareUnicodeCharacter{2081}{\textsubscript{1}}
\DeclareUnicodeCharacter{2082}{\textsubscript{2}}
\DeclareUnicodeCharacter{2083}{\textsubscript{3}}
\DeclareUnicodeCharacter{2084}{\textsubscript{4}}
\DeclareUnicodeCharacter{2085}{\textsubscript{5}}
\DeclareUnicodeCharacter{2086}{\textsubscript{6}}
\DeclareUnicodeCharacter{2087}{\textsubscript{7}}
\DeclareUnicodeCharacter{2088}{\textsubscript{8}}
\DeclareUnicodeCharacter{2089}{\textsubscript{9}}
\DeclareUnicodeCharacter{00B0}{\textdegree}

% Counter behaviour
\usepackage{chngcntr}
\counterwithout{figure}{section}
\counterwithout{table}{section}
\counterwithout{equation}{section}

% Caption formatting
\usepackage[font=small,labelfont=bf,textfont=it]{caption}
\captionsetup{justification=raggedright,singlelinecheck=false}
\captionsetup[figure]{position=bottom,skip=10pt}
\captionsetup[table]{position=top,skip=10pt}
\AtBeginEnvironment{figure}{\raggedright}

% Titlesec formatting
\usepackage{titlesec}
\titlespacing{\section}{0pt}{12pt}{6pt}
\titlespacing{\subsection}{0pt}{10pt}{4pt}
\titlespacing{\subsubsection}{0pt}{8pt}{3pt}
\titleformat{\section}{\large\bfseries}{\thesection}{1em}{}
\titleformat{\subsection}{\normalsize\bfseries}{\thesubsection}{1em}{}
\titleformat{\subsubsection}{\normalsize\itshape}{\thesubsubsection}{1em}{}

% Ensure each top-level section starts on a new page
\let\oldsection\section
\renewcommand{\section}{\clearpage\oldsection}

% Graphics support for figures
\usepackage{graphicx}

% Small visual tweaks
\usepackage{xcolor}
\definecolor{codebg}{RGB}{245,245,245}
\definecolor{codeborder}{RGB}{200,200,200}
\definecolor{codefg}{RGB}{50,50,50}

% Setup listings safely (don't override global styles)
\usepackage{listings}
\lstset{backgroundcolor=\color{codebg},basicstyle=\small\ttfamily\color{codefg},frame=single,framerule=0.5pt,framesep=5pt}

% Keep title metadata to be inserted by build system
\ifLuaTeX
  \usepackage{selnolig}  % disable illegal ligatures
\fi
\IfFileExists{bookmark.sty}{\usepackage{bookmark}}{\usepackage{hyperref}}
\IfFileExists{xurl.sty}{\usepackage{xurl}}{} % add URL line breaks if available
\urlstyle{same}
\hypersetup{
  colorlinks=true,
  linkcolor={blue},
  filecolor={blue},
  citecolor={blue},
  urlcolor={blue},
  pdfcreator={LaTeX via pandoc}}

\author{}
\date{}

\begin{document}

{
\hypersetup{linkcolor=black}
\setcounter{tocdepth}{3}
\tableofcontents
}
\setstretch{1.0}
\section{Appendix D: Neural Encoding Efficiency on
Time-Series}\label{sec:app_neural_encoding}

\subsection{Objective}\label{objective}

Comprehensive neural encoding analysis including spike‑train generation,
temporal dynamics, population coding, mutual information, and adaptation
mechanisms for olfactory receptor neurons.

\subsection{Methods (src)}\label{methods-src}

\begin{itemize}
\tightlist
\item
  \passthrough{\lstinline!src/case\_studies/neural\_encoding.py!}

  \begin{itemize}
  \tightlist
  \item
    \passthrough{\lstinline!generate\_spike\_trains(stimuli, dt, baseline\_rate, max\_rate, dynamics)!}
    --- realistic spike generation
  \item
    \passthrough{\lstinline!analyze\_spike\_train\_statistics(spike\_data)!}
    --- ISI, CV, Fano factor
  \item
    \passthrough{\lstinline!temporal\_coding\_analysis(spike\_data, stimulus\_times)!}
    --- latency and precision metrics
  \item
    \passthrough{\lstinline!population\_coding\_analysis(population\_responses, labels)!}
    --- PCA, LDA, correlation structure
  \item
    \passthrough{\lstinline!mutual\_information\_analysis(responses, stimuli)!}
    --- information‑theoretic metrics
  \item
    \passthrough{\lstinline!odor\_discrimination\_analysis(responses, odor\_ids, time\_windows)!}
    --- discrimination performance
  \item
    \passthrough{\lstinline!adaptation\_dynamics\_analysis(spike\_data, stimulus\_duration)!}
    --- adaptation characterization
  \item
    \passthrough{\lstinline!information\_rate\_time\_series(responses, dt\_s, noise\_std)!}
    --- channel‑capacity estimation
  \item
    \passthrough{\lstinline!rate\_coding\_metrics(responses, labels)!}
    --- separability and discriminability metrics
  \end{itemize}
\end{itemize}

\subsection{Script and outputs}\label{script-and-outputs}

\begin{itemize}
\tightlist
\item
  Script:
  \passthrough{\lstinline!scripts/generate\_neural\_encoding\_analysis.py!}
\item
  Data:
  \passthrough{\lstinline!output/data/neural\_encoding\_comprehensive.npz!}
\item
  Figure:
  \passthrough{\lstinline!figures/neural\_encoding\_comprehensive\_analysis.png!}
\end{itemize}

\subsection{Figure}\label{figure}

\begin{figure}[h]
\centering
\includegraphics[width=0.9\textwidth]{figures/neural_encoding_comprehensive_analysis.png}
\caption{Neural encoding analyses generated deterministically by `scripts/generate_neural_encoding_analysis.py` using `src/case_studies/neural_encoding.py`. Panels include spike trains, temporal precision, population PCA, ISI statistics, and mutual information metrics.}
\label{fig:app_neural_encoding_full}
\end{figure}

\subsection{Equation references}\label{equation-references}

\begin{itemize}
\tightlist
\item
  Information rate: see \eqref{eq:channel_capacity_gloss}
\item
  Response time model: see \eqref{eq:response_time_components}
\end{itemize}

\subsection{Reproducibility}\label{reproducibility}

\begin{itemize}
\tightlist
\item
  Run:
  \passthrough{\lstinline!python3 scripts/generate\_neural\_encoding\_analysis.py!}
\item
  Artifacts saved to \passthrough{\lstinline!output/data/!} and
  \passthrough{\lstinline!figures/!}.
\item
  Deterministic seeds:
  \passthrough{\lstinline!src/config.set\_random\_seed(42)!} for
  surrogate time‑series.
\end{itemize}

\subsection{Cross‑references}\label{crossreferences}

\begin{itemize}
\tightlist
\item
  Methods: \Cref{sec:methodology}
\item
  Symbols: \Cref{sec:symbols_glossary}
\item
  Math appendix: \Cref{sec:mathematical_appendix}
\end{itemize}

\end{document}
